\section{Limitations and Threats to Validity}
\label{sec:limitations}

\paragraph{Construct validity.}
Controlled labels are construction labels written by the system's authors; they
encode what scenarios were built to contain. The external outcome (SLA miss)
measures timeliness, not supervisory execution quality, and investigation and
escalation are proxies there. Risk scores have no validated meaning as
probabilities or severities.

\paragraph{Internal validity.}
Populations for prioritization and ablation come from SAT-SA's own generator,
which injects the behaviours its detectors target, favouring SAT-SA. Precision
depends on database layout (shared versus isolated). Detector thresholds were set
by design, not fitted; we made no changes on evaluation data, but they were also
never tuned on independent development data.

\paragraph{External validity.}
There is no SOC data from a real organisation. The external dataset is one IT
service desk. Generated populations have \V{PR01-entities} entities, making top-$k$
estimates coarse. The ablation population analysis uses five seeds and is
exploratory.

\paragraph{Statistical conclusion validity.}
Intervals are descriptive, seeded and uncorrected across many comparisons; no
hypothesis test is made. Orchestration measurements are repeated runs on one
machine and do not generalise across hardware, databases or load. Recovery was
tested with injected exceptions at worker-method boundaries, not real process
kills, partitions or database failover.

\paragraph{Human factors.}
The human evaluation protocol is ready but \textbf{the study was not executed};
there are no participants and no usability or decision-quality results. The expert
label pipeline is ready, but \textbf{no expert-labelled data exist}.

\paragraph{Security.}
Integrity results cover \V{T01-expected} tested mutations on one workflow; key
compromise, key management, insider re-signing and false submissions are out of
scope. Tenant isolation was not tested adversarially.

\paragraph{Deployment.}
All measurements used SQLite and local storage on one machine. The local API and
worker processes passed \V{D01-pass} HTTP smoke checks (\V{D01-fail} failed,
\V{D01-notrun} not run). Docker was not installed, so container images and the
compose topology were not executed; PostgreSQL was not live-tested (no connection
string was available); SeaweedFS object storage and the CI topology job were not
executed; no hosted deployment exists (Table~\ref{tab:deployment}). No claim of
production readiness or scale is made.

% Generated by scripts/build_paper_assets.py from satsa-evidence-freeze-v2; do not edit.
% Sources: O01=EXP-O01-orchestration-overhead (manifest d758a2801f158cc2)
\begin{table}[t]
\centering\small
\caption{Deployment verification status and runtime characteristics.}
\label{tab:deployment}
\begin{adjustbox}{max width=\linewidth}
\begin{tabular}{ll}
\toprule
Layer & Status \\
\midrule
API + worker processes (SQLite, local storage) & local HTTP smoke \V{D01-pass}/18 checks passed \\
Container images and compose topology & configured; not executed (Docker unavailable) \\
PostgreSQL & implemented; not live-tested \\
S3-compatible storage (SeaweedFS) & configured; not executed \\
CI topology job & configured; not executed \\
Hosted deployment & not executed \\
Workflow processing (O01, direct) & median \V{O01-direct-median} s; finalization \V{O01-final-median} s; verification \V{O01-verify-median} s \\
\bottomrule
\end{tabular}
\end{adjustbox}
\end{table}

